% !TEX root = ../main.tex
\section{Semiconductor Foundations}
\label{sec:semiconductor_foundations}

Before the working principle of a laser diode can be discussed in detail, preliminary
knowledge about carrier and photon statistics is required. Furthermore, the diode
current-voltage characteristics are discussed in this section, as they are not the
main focus of this laboratory course.

\subsection{Carrier Statistics}
\label{sec:carrier_statistics}

In thermal equilibrium, the probability that an electronic state of energy $E$ is
occupied is given by the Fermi-Dirac distribution, with the Fermi level $E_{\mathrm{F}}$,
the Boltzmann constant $k_{\mathrm{B}}$ and the temperature $T$:
\begin{equation}
	f(E)=\frac{1}{1+\exp\left[(E-E_{\mathrm{F}})/k_{\mathrm{B}}T\right]}.
\end{equation}

However, under non-equilibrium conditions (e.g., under an applied external voltage),
electron and hole densities can take arbitrary values in principle. In many
practically relevant cases, the carriers within each band are locally in thermal equilibrium
among themselves, but not with other bands. Instead of a constant Fermi level,
electron and hole densities are described locally at position $\mathbf{r}$ using quasi-Fermi levels
$F_{\mathrm{n}}(\mathbf{r})$ and $F_{\mathrm{p}}(\mathbf{r})$, and the respective
occupation probabilities are given by:
\begin{equation}
	f_{\mathrm{n}}(E, \mathbf{r}) = \frac{1}{1+\exp\left[(E-F_{\mathrm{n}}(\mathbf{r}))/k_{\mathrm{B}}T\right]}, \qquad
	f_{\mathrm{p}}(E, \mathbf{r}) = \frac{1}{1+\exp\left[(F_{\mathrm{p}}(\mathbf{r})-E)/k_{\mathrm{B}}T\right]}.
	\label{eq:occupation_probabilities}
\end{equation}
The electron density $n(\mathbf{r})$ and the hole density $p(\mathbf{r})$ are obtained
by integrating the product of the occupation probability and the density of states
$D_{\mathrm{n}}(E)$ or $D_{\mathrm{p}}(E)$ of the electrons or holes, respectively.
The integration runs over the conduction band above its edge $E_{\mathrm{C}}$ and
over the valence band below its edge $E_{\mathrm{V}}$:
\begin{equation}
	n(\mathbf{r}) = \int_{E_{\mathrm{C}}}^{\infty}D_{\mathrm{n}}(E)\, f_{\mathrm{n}}
	(E,\mathbf{r}) \,\mathrm{d}E, \qquad p(\mathbf{r}) = \int_{-\infty}^{E_{\mathrm{V}}}
	D_{\mathrm{p}}(E)\, f_{\mathrm{p}}(E,\mathbf{r}) \,\mathrm{d}E.
\end{equation}
Using the Boltzmann approximation, the integrals can be evaluated analytically, resulting
in the following expressions with the intrinsic carrier density $n_{\mathrm{i}}$ and the
intrinsic Fermi level $E_{\mathrm{i}}(\mathbf{r})$:
\begin{equation}
	n(\mathbf{r}) = n_{\mathrm{i}}\exp\left[\frac{F_{\mathrm{n}}(\mathbf{r}) - E_{\mathrm{i}}(\mathbf{r})}{k_{\mathrm{B}}T}\right], \qquad
	p(\mathbf{r}) = n_{\mathrm{i}}\exp\left[\frac{E_{\mathrm{i}}(\mathbf{r}) - F_{\mathrm{p}}(\mathbf{r})}{k_{\mathrm{B}}T}\right].
\end{equation}

The splitting $\Delta F(\mathbf{r})=F_{\mathrm{n}}(\mathbf{r})-F_{\mathrm{p}}(\mathbf{r}
)$ is a local quantity: it fixes the carrier product $np$ at each position. The applied
voltage $V$ is a global quantity: it fixes the difference between the quasi-Fermi
levels of the two contacts, with $V>0$ for forward bias and the elementary charge $e$:
\begin{equation}
	eV=F_{\mathrm{n}}(\text{n-contact})-F_{\mathrm{p}}(\text{p-contact}).
\end{equation}

\subsection{Joint Density of States}
\label{sec:band_structure}

Near the zone center of a direct semiconductor, the conduction and valence bands
can often be approximated by parabolic functions:
\begin{equation}
	E_{\mathrm{c}}(k)=E_{\mathrm{C}}+\frac{\hbar^{2}k^{2}}{2m_{\mathrm{e}}^{*}},\qquad
	E_{\mathrm{v}}(k)=E_{\mathrm{V}}-\frac{\hbar^{2}k^{2}}{2m_{\mathrm{h}}^{*}},\qquad
	E_{\mathrm{g}}=E_{\mathrm{C}}-E_{\mathrm{V}}.
	\label{eq:band_dispersion}
\end{equation}
Here, $k$ is the magnitude of the wavevector, $\hbar$ the reduced Planck constant,
$m_{\mathrm{e}}^{*}$ and $m_{\mathrm{h}}^{*}$ the effective masses of electrons and
holes, $E_{\mathrm{C}}$ and $E_{\mathrm{V}}$ the conduction band minimum and the
valence band maximum, and $E_{\mathrm{g}}$ the band gap. The energy difference
$E_{\mathrm{cv}}(k)$ between the two bands can be stated using the reduced mass
$m_{\mathrm{r}}$:
\begin{equation}
	E_{\mathrm{cv}}(k)=E_{\mathrm{c}}(k)-E_{\mathrm{v}}(k)=E_{\mathrm{g}}+\frac{\hbar^{2}k^{2}}{2m_{\mathrm{r}}}
	,\qquad \frac{1}{m_{\mathrm{r}}}=\frac{1}{m_{\mathrm{e}}^{*}}+\frac{1}{m_{\mathrm{h}}^{*}}
	.
	\label{eq:parabolic_bands}
\end{equation}

An optical transition in a semiconductor conserves energy as well as momentum of the
participating electron and photon. As the photon wavevector is small compared to
the extent of the Brillouin Zone, transitions can be approximated as vertical in $\mathbf{k}$-space.

The rate of optical transitions at energy $E$ depends on how many pairs of electronic
states are separated by exactly this energy. The number of available transitions
per unit volume and per energy is called the joint density of states
$D_{\mathrm{j}}(E)$ and is defined as:
\begin{equation}
	\text{\setstretch{1}
	\begin{tabular}{@{}c@{}}
		number of transitions per unit volume \\
		with energy in $[E_{0}, E_{1}]$
	\end{tabular}}= \int_{E_{0}}^{E_{1}}D_{\mathrm{j}}(E) \, \mathrm{d}E.
\end{equation}
Analogous to the electronic density of states, $D_{\mathrm{j}}$ can be defined
by counting the $\mathbf{k}$-states on the surface $S(\hbar\omega)$ of constant
transition energy $E_{\mathrm{cv}}(\mathbf{k})=\hbar\omega$, where $\omega$ is the
angular frequency of the photon:
\begin{equation}
	D_{\mathrm{j}}(\hbar \omega)= \frac{2}{V}\int_{S(\hbar \omega)}\frac{\mathrm{d}^{2}S}{\left(
	\frac{2\pi}{L}\right)^{3}\lvert \nabla_{\mathbf{k}}E_{\mathrm{cv}}(\mathbf{k}) \rvert
	}.
	\label{eq:joint_dos_surface}
\end{equation}
The factor 2 counts the two spin orientations, and $V=L^{3}$ is the volume of the
periodic boundary box with edge length $L$. By using the parabolic band approximation,
this integral can be solved analytically, with the Heaviside step function $\Theta$:
\begin{equation}
	D_{\mathrm{j}}(\hbar\omega)=\frac{1}{2\pi^{2}}\left(\frac{2m_{\mathrm{r}}}{\hbar^{2}}
	\right)^{3/2}\sqrt{\hbar\omega-E_{\mathrm{g}}}\;\Theta(\hbar\omega-E_{\mathrm{g}}
	).
\end{equation}

\subsection{Photon Mode Density and Statistics}
\label{sec:photons}

The allowed states of a photon field can be constructed using periodic boundary conditions
for the photon wave $\psi_{\mathbf{k}}$. These conditions result in a finite
set of allowed $\mathbf{k}$ vectors, with integers $l$, $m$, $n$ and the edge length
$L$ of the periodic boundary box:
\begin{equation}
	k_{x}=l \frac{2\pi}{L}, \quad k_{y}=m \frac{2\pi}{L}, \quad k_{z}=n \frac{2\pi}{L}
	.
\end{equation}

Analogous to the electronic density of states, the photon mode density of states $N
(E)$ can be constructed by counting the number of available states with energy $E$.
As the periodic boundary conditions lead to the same density of states in
$\mathbf{k}$-space as for the electrons, \cref{eq:joint_dos_surface} can be used
analogously, with a polarization factor of 2 (two transverse polarization states)
instead of the spin factor of 2. Substituting the photon energy relation
$E(\mathbf{k})=\hbar \lvert \mathbf{k}\rvert c /n_{\mathrm{r}}$ with the vacuum speed
of light $c$ for a dispersion-free medium, i.e. a constant refractive index
$n_{\mathrm{r}}$, the integral simplifies to:
\begin{equation}
	N(E)=\frac{n_{\mathrm{r}}^{3}E^{2}}{\pi^{2}\hbar^{3}c^{3}}.
\end{equation}

The spectral photon density $S(E)$, i.e. the number of photons per unit volume and
unit energy at energy $E$, is the product of the mode density $N(E)$ and the average
photon occupation number per mode $n_{\mathrm{ph}}(E)$:
\begin{equation}
	S(E)=N(E) \cdot n_{\mathrm{ph}}(E).
\end{equation}

As photons are bosonic particles, they obey the Bose-Einstein Statistics. In thermal
equilibrium with the lattice at temperature $T$, the occupation number $n_{\mathrm{ph}}$
per mode with energy $E$ is therefore given by:
\begin{equation}
	n_{\mathrm{ph}}(E)=\frac{1}{\exp\left( \dfrac{E}{k_{\mathrm{B}}T} \right) - 1}.
\end{equation}

\subsection{Diode Current-Voltage Characteristics}
\label{sec:carrier_injection}

The current-voltage characteristic of an ideal diode is described by the
Shockley equation for the current $I$, with the saturation current $I_{0}$ and the
voltage $V_{\mathrm{a}}$ across the junction:
\begin{equation}
	I = I_{0}\left[ \exp\left(\frac{eV_{\mathrm{a}}}{k_{\mathrm{B}}T}\right) - 1 \right
	].
\end{equation}
In the forward direction ($V_{\mathrm{a}}> 0$), the current increases
exponentially with the voltage, while in the reverse direction ($V_{\mathrm{a}}< 0$)
it saturates at the small reverse current $-I_{0}$. Real diodes are often
described by replacing $k_{\mathrm{B}}T$ with $n_{\mathrm{id}}k_{\mathrm{B}}T$, where
the ideality factor $n_{\mathrm{id}}$ usually lies between 1 and 2.

In a real diode, part of the applied voltage $V$ drops across the series resistance
$R_{\mathrm{s}}$ of the semiconductor bulk and the contacts. Only the remaining
voltage $V_{\mathrm{a}}= V - IR_{\mathrm{s}}$ is present at the junction, so the characteristic
becomes:
\begin{equation}
	I = I_{0}\left[\exp\left(\frac{e(V - IR_{\mathrm{s}})}{n_{\mathrm{id}} k_{\mathrm{B}}T}\right) -
	1\right].
	\label{eq:schockley_real}
\end{equation}